\section{模拟设置}

本研究使用由 CLS 合作组以周期性边界条件产生的、含 2+1 味 clover 费米子与树级 Symanzik 规范作用的组态~\cite{Bruno:2014jqa}。详细参数列于表~\ref{table:cls}。注意价夸克采用的是 $m_\pi=547$ MeV 而非 333 MeV，以获得更好的信号。物理上，当增强因子 $\gamma$ 足够大时，软函数将与介子质量无关。

为计算式~(\ref{eq:FF}) 中的形状因子，我们生成墙源传播子
\begin{align}
 S_w(x,t,t';\vec{p})=\sum_{\vec{y}} S(t,\vec{x};t',\vec{y})e^{i\vec{p}\cdot(\vec{y}-\vec{x})},
\end{align}
在固定到 Coulomb 规范的组态上，对初、末介子态均在 $t'=0$ 与 $t_{\rm sep}$ 处取墙源。$S$ 是从 $(t',\vec{y})$ 到 $(t,\vec{x})$ 的夸克传播子。于是可以构造与式~(\ref{eq:FF}) 中形状因子对应的三点函数（3pt）
\begin{align}
&C_3(b_\perp,P^z;p^z,t_{\rm sep},t)\\
&=\frac{1}{L^3} \sum_x\textrm{Tr} \langle  S_{w}^\dagger(\vec{x}+\vec{b},t,0;-\vec{p})\gamma_5\Gamma S_{w}(\vec{x}+\vec{b},t,t_{\rm sep};\vec{p}) \nonumber\\
&\;\;\quad\quad \quad \times S_{w}^\dagger(\vec{x},t,t_{\rm sep};-\vec{P}+\vec{p}) \gamma_{5}\Gamma
 S_{w}(\vec{x},t,0;\vec{P}-\vec{p})\rangle\nonumber.
\end{align}
夸克动量 $\vec{p}=(\vec 0_{\perp},p^z)$，并对反夸克传播子使用了关系 $\gamma_5 S^{\dagger}(x,y)\gamma_5=S(y,x)$。我们试验了若干种 $\Gamma$ 的选取，最终采用单位 Dirac 矩阵 $\Gamma=I$，因为它信号最好，且在大 $P^z$ 极限下描述领头扭度的光锥贡献。注意 $\Gamma=\gamma_4$ 情形在大 $P^z$ 极限下是次领头的，尽管其激发态污染可能更小，因而不太适用。

通过在全部 48 个时间片上以夸克动量 ${p}^z=(-2,-1,0,1,2) \times 2\pi/(La)$ 生成墙源传播子，我们可以对任意 $t$ 与 $t_{\rm sep}$、介子动量 $P_z$ 从 0 到 $8\pi/(La)$（$\sim 2.1$ GeV）最大限度地提高三点函数的统计量。利用带一个激发态的标准三点函数参数化，$C_3(b_\perp,P^z,t_{\rm sep},t)$ 与裸 $F(b_\perp,P^z)$ 的关系为
\begin{align}
& C_3(b_\perp,P^z;p^z,t_{\rm sep},t)=\frac{A_w(p_z)^2}{(2E)^2}e^{-Et_{\rm sep}}\big[F(b_\perp,P^z)\nonumber\\
&\quad\quad+c_1(e^{-\Delta E t}+e^{-\Delta E (t_{\rm sep}-t)})+c_2e^{-\Delta E t_{\rm sep}}\big].
 \label{eq:3pt}
\end{align}
$A_w$ 是 Coulomb 规范固定墙源（CFW）π介子插值场的矩阵元，$E=\sqrt{m_{\pi}^2+{P}^{z2}}$ 是π介子能量，$\Delta E$ 是π介子与其第一激发态之间的能隙，$c_{1,2}$ 是激发态污染参数。注意 $A_w^2$ 的 $p_z$ 依赖因子将被抵消{。}

同样的墙源传播子可用于计算与裸准 TMDWF 相联系的二点函数，
\begin{align}
&C_2(b_\perp,P^z;p_z,\ell,t)=\frac{1}{L^3\sqrt{Z_E(2\ell,b_\perp)}}\sum_x \textrm{Tr}e^{i\vec{P}\cdot \vec{x}}\nonumber\\
&\times\langle  S_w^\dagger(\vec{x}+\vec{b},t,0;-\vec{p}){\cal W}(\vec b,\ell) \gamma_5\Gamma_\Phi  S_w(\vec{x},t,0;P^z-\vec{p}) \rangle\nonumber\\
& =\frac{A_w(p_z)A_p}{2E}e^{-Et}\phi_\ell(0,b_\perp,P^z,\ell)(1+c_0 e^{-\Delta E t}),\label{eq:2pt}
\end{align}
其中同样用一个激发态做参数化混合。
$A_p$ 是点汇π介子插值场的矩阵元。当我们用 $\phi_\ell(0,0,P^z,0)$ 归一化 $\phi_\ell(0,b_\perp,P^z,\ell)$ 时它会被去除。我们选取 $\Gamma_\Phi=\gamma^t\gamma_5$ 来定义式~(\ref{eq:phi_def}) 中的波函数振幅。基于文献~\cite{Shanahan:2019zcq,Shanahan:2020zxr}中采用类似钉形规范链算符的准 TMDPDF 研究，混合效应{在把各种贡献求和后可能相当可观。在补充材料中，我们报告了使用 A654 系综的类似模拟。我们发现对横向间隔 $b_\perp\sim 0.6{\rm fm}$，混合效应可达 $5\%$ 量级。这些效应将作为系统误差之一纳入以下分析，而关于混合效应的全面研究将在未来进行。}


\begin{figure}[!th]
\includegraphics[width=0.45\textwidth]{images/02_tmdwf_mod_with_L_P3_b3_t6.pdf}
\caption{ 以 $\{P^z, b_{\perp}, t\}=\{6\pi/L, 3a, 6a\}$ 为例，$z=0$ 的准 TMDWF 随 $\ell$ 的变化结果，以及用于减除的 Wilson 圈的平方根。所有结果均以其 $\ell=0$ 处的值归一。}\label{fig:L_dependence}
\end{figure}

π介子态的色散关系、测量直方图的统计检验以及组态间自关联的信息见补充材料~\cite{supplemental}。
